% GENERATED FILE. Edit canonical lessons, then run npm run generate:books.
% canonical-source-hash: fnv1a64:5fef52fdf346854f
% canonical-lessons: JA-C487-shoubai, JA-C487-shounin, JA-C487-tenchou, JA-C487-uriage, JA-C487-juugyouin

\chapter{Running a Business}
\label{ch:ja-running-a-business}
\hlchaptermodality{487}

\begin{chapteropening}
\textbf{By the end of this chapter:} \emph{I can talk about a shop and the people who work in it.}

Complete the last lesson of chapter 487: I can talk about a shop and the people who work in it.
\end{chapteropening}

\section[\texorpdfstring{\ja{しょうばい} (shōbai)}{しょうばい (shōbai)}]{\ja{しょうばい} (shōbai) — business; trade}

\label{lesson:JA-C487-shoubai}



\begin{quote}
Before the new one: say the Japanese for to fit in, then the Japanese for to peel off.
\end{quote}

\subsection*{You'll want to know: \ja{しょうばい}}
\textbf{\ja{しょうばい}} — \emph{shōbai} — \textquotedblleft{}business; trade\textquotedblright{}.

\begin{quote}
\textbf{\ja{そふは} \ja{しょうばいを} \ja{しています。}} — \emph{sofu wa shōbai o shiteimasu} — \textquotedblleft{}My grandfather runs a business.\textquotedblright{}
\end{quote}

\begin{cousinweb}
\textbf{\ja{しょうばい}} (\emph{shōbai}) is Sino-Japanese (kango): two Chinese characters read with their on'yomi, readings borrowed from Chinese. They are \emph{shō}, \textquotedblleft{}selling\textquotedblright{} (Middle Chinese \emph{syang}) and \emph{bai}, \textquotedblleft{}sell\textquotedblright{} (\emph{meaH}) (KANJIDIC2; Middle Chinese in Baxter's transcription).
\end{cousinweb}

\begin{grammarlens}[title={What you've built}]
One more word for describing what you lived and felt.
\end{grammarlens}

\subsection*{Your turn}
\begin{itemize}
\raggedright
  \item \emph{Say it:} \emph{shōbai}
  \item \emph{Say it:} \emph{shōbai}, once more
  \item \emph{Say it:} \emph{sofu wa shōbai o shiteimasu} — My grandfather runs a business.
  \item \emph{Recall:} say the Japanese for to fit in, then the Japanese for to peel off, then say \emph{shōbai} again
\end{itemize}

\begin{tcolorbox}[breakable,colback=teal!4,colframe=teal!35!black,title={Before you move on}]
What does \ja{しょうばい} mean? (\textquotedblleft{}Business; trade\textquotedblright{}.) Say it once more.
\end{tcolorbox}
\section[\texorpdfstring{\ja{しょうにん} (shōnin)}{しょうにん (shōnin)}]{\ja{しょうにん} (shōnin) — a merchant; a trader}

\label{lesson:JA-C487-shounin}



\begin{quote}
Before the new one: say the Japanese for to peel off, then the Japanese for business.
\end{quote}

\subsection*{You'll want to know: \ja{しょうにん}}
\textbf{\ja{しょうにん}} — \emph{shōnin} — \textquotedblleft{}a merchant; a trader\textquotedblright{}.

\begin{quote}
\textbf{\ja{しょうにんが} \ja{いちばで} \ja{やさいを} \ja{うります。}} — \emph{shōnin ga ichiba de yasai o urimasu} — \textquotedblleft{}Traders sell vegetables at the market.\textquotedblright{}
\end{quote}

\begin{cousinweb}
\textbf{\ja{しょうにん}} (\emph{shōnin}) is Sino-Japanese (kango): two Chinese characters read with their on'yomi, readings borrowed from Chinese. They are \emph{shō}, \textquotedblleft{}merchant\textquotedblright{} (Middle Chinese \emph{syang}) and \emph{nin}, \textquotedblleft{}person\textquotedblright{} (\emph{nyin}) (KANJIDIC2; Middle Chinese in Baxter's transcription).
\end{cousinweb}

\begin{grammarlens}[title={What you've built}]
One more word for describing what you lived and felt.
\end{grammarlens}

\subsection*{Your turn}
\begin{itemize}
\raggedright
  \item \emph{Say it:} \emph{shōnin}
  \item \emph{Say it:} \emph{shōnin}, once more
  \item \emph{Say it:} \emph{shōnin ga ichiba de yasai o urimasu} — Traders sell vegetables at the market.
  \item \emph{Recall:} say the Japanese for to peel off, then the Japanese for business, then say \emph{shōnin} again
\end{itemize}

\begin{tcolorbox}[breakable,colback=teal!4,colframe=teal!35!black,title={Before you move on}]
What does \ja{しょうにん} mean? (\textquotedblleft{}A merchant; a trader\textquotedblright{}.) Say it once more.
\end{tcolorbox}
\section[\texorpdfstring{\ja{てんちょう} (tenchō)}{てんちょう (tenchō)}]{\ja{てんちょう} (tenchō) — a shop manager}

\label{lesson:JA-C487-tenchou}



\begin{quote}
Before the new one: say the Japanese for business, then the Japanese for a merchant.
\end{quote}

\subsection*{You'll want to know: \ja{てんちょう}}
\textbf{\ja{てんちょう}} — \emph{tenchō} — \textquotedblleft{}a shop manager\textquotedblright{}.

\begin{quote}
\textbf{\ja{てんちょうは} \ja{まいあさ} \ja{みせを} \ja{あけます。}} — \emph{tenchō wa maiasa mise o akemasu} — \textquotedblleft{}The manager opens the shop every morning.\textquotedblright{}
\end{quote}

\begin{cousinweb}
\textbf{\ja{てんちょう}} (\emph{tenchō}) is Sino-Japanese (kango): two Chinese characters read with their on'yomi, readings borrowed from Chinese. They are \emph{ten}, \textquotedblleft{}store\textquotedblright{} (Middle Chinese \emph{temH}) and \emph{chō}, \textquotedblleft{}leader\textquotedblright{} (\emph{trjangX}) (KANJIDIC2; Middle Chinese in Baxter's transcription).
\end{cousinweb}

\begin{grammarlens}[title={What you've built}]
One more word for describing what you lived and felt.
\end{grammarlens}

\subsection*{Your turn}
\begin{itemize}
\raggedright
  \item \emph{Say it:} \emph{tenchō}
  \item \emph{Say it:} \emph{tenchō}, once more
  \item \emph{Say it:} \emph{tenchō wa maiasa mise o akemasu} — The manager opens the shop every morning.
  \item \emph{Recall:} say the Japanese for business, then the Japanese for a merchant, then say \emph{tenchō} again
\end{itemize}

\begin{tcolorbox}[breakable,colback=teal!4,colframe=teal!35!black,title={Before you move on}]
What does \ja{てんちょう} mean? (\textquotedblleft{}A shop manager\textquotedblright{}.) Say it once more.
\end{tcolorbox}
\section[\texorpdfstring{\ja{うりあげ} (uriage)}{うりあげ (uriage)}]{\ja{うりあげ} (uriage) — sales; takings}

\label{lesson:JA-C487-uriage}



\begin{quote}
Before the new one: say the Japanese for a merchant, then the Japanese for a shop manager.
\end{quote}

\subsection*{You'll want to know: \ja{うりあげ}}
\textbf{\ja{うりあげ}} — \emph{uriage} — \textquotedblleft{}sales; takings\textquotedblright{}.

\begin{quote}
\textbf{\ja{せんしゅうの} \ja{うりあげは} \ja{よかったです。}} — \emph{senshū no uriage wa yokatta desu} — \textquotedblleft{}Last week's sales were good.\textquotedblright{}
\end{quote}

\begin{cousinweb}
\textbf{\ja{うりあげ}} (\emph{uriage}) is a native Japanese compound (wago), inherited rather than borrowed from Chinese: built from \emph{uri}, \textquotedblleft{}selling\textquotedblright{} and \emph{age}, \textquotedblleft{}raising\textquotedblright{} (JMdict, KANJIDIC2).
\end{cousinweb}

\begin{grammarlens}[title={What you've built}]
One more word for describing what you lived and felt.
\end{grammarlens}

\subsection*{Your turn}
\begin{itemize}
\raggedright
  \item \emph{Say it:} \emph{uriage}
  \item \emph{Say it:} \emph{uriage}, once more
  \item \emph{Say it:} \emph{senshū no uriage wa yokatta desu} — Last week's sales were good.
  \item \emph{Recall:} say the Japanese for a merchant, then the Japanese for a shop manager, then say \emph{uriage} again
\end{itemize}

\begin{tcolorbox}[breakable,colback=teal!4,colframe=teal!35!black,title={Before you move on}]
What does \ja{うりあげ} mean? (\textquotedblleft{}Sales; takings\textquotedblright{}.) Say it once more.
\end{tcolorbox}
\section[\texorpdfstring{\ja{じゅうぎょういん} (jūgyōin)}{じゅうぎょういん (jūgyōin)}]{\ja{じゅうぎょういん} (jūgyōin) — an employee; a worker}

\label{lesson:JA-C487-juugyouin}



\begin{quote}
Before the new one: say the Japanese for a shop manager, then the Japanese for sales.
\end{quote}

\subsection*{You'll want to know: \ja{じゅうぎょういん}}
\textbf{\ja{じゅうぎょういん}} — \emph{jūgyōin} — \textquotedblleft{}an employee; a worker\textquotedblright{}.

\begin{quote}
\textbf{\ja{この} \ja{みせには} \ja{じゅうぎょういんが} \ja{ごにん} \ja{います。}} — \emph{kono mise ni wa jūgyōin ga gonin imasu} — \textquotedblleft{}This shop has five employees.\textquotedblright{}
\end{quote}

\begin{cousinweb}
\textbf{\ja{じゅうぎょういん}} (\emph{jūgyōin}) is Sino-Japanese (kango): three Chinese characters read with their on'yomi, readings borrowed from Chinese. They are \emph{jū}, \textquotedblleft{}accompany\textquotedblright{} (Middle Chinese \emph{dzjowng}), \emph{gyō}, \textquotedblleft{}business\textquotedblright{} (\emph{ngjaep}) and \emph{in}, \textquotedblleft{}employee\textquotedblright{} (\emph{hjwen}) (KANJIDIC2; Middle Chinese in Baxter's transcription).
\end{cousinweb}

\begin{grammarlens}[title={What you've built}]
One more word for describing what you lived and felt.
\end{grammarlens}

\subsection*{Your turn}
\begin{itemize}
\raggedright
  \item \emph{Say it:} \emph{jūgyōin}
  \item \emph{Say it:} \emph{jūgyōin}, once more
  \item \emph{Say it:} \emph{kono mise ni wa jūgyōin ga gonin imasu} — This shop has five employees.
  \item \emph{Recall:} say the Japanese for a shop manager, then the Japanese for sales, then say \emph{jūgyōin} again
\end{itemize}

\begin{tcolorbox}[breakable,colback=teal!4,colframe=teal!35!black,title={Before you move on}]
What does \ja{じゅうぎょういん} mean? (\textquotedblleft{}An employee; a worker\textquotedblright{}.) Say it once more.
\end{tcolorbox}
